\documentclass[conference]{IEEEtran}
\IEEEoverridecommandlockouts
\usepackage{cite}
\usepackage{amsmath,amssymb,amsfonts}
\usepackage{algorithmic}
\usepackage{graphicx}
\usepackage{textcomp}
\usepackage{xcolor}
\usepackage{booktabs}
\usepackage{algorithm}

\def\BibTeX{{\rm B\kern-.05em{\sc i\kern-.025em b}\kern-.08em
    T\kern-.1667em\lower.7ex\hbox{E}\kern-.125emX}}

\begin{document}

\title{Confidence-Gated Selective Consultation in LLM-Based Multi-Agent Decision Systems}

\author{
\IEEEauthorblockN{Palak Malpani}
\IEEEauthorblockA{\textit{School of Computer Science and Engineering (SCOPE)} \\
\textit{Vellore Institute of Technology, Chennai}\\
Chennai, Tamil Nadu, India \\
palak.malpani2024@vitstudent.ac.in}
\and
\IEEEauthorblockN{Ramakrishnan P. H.}
\IEEEauthorblockA{\textit{School of Computer Science and Engineering (SCOPE)} \\
\textit{Vellore Institute of Technology, Chennai}\\
Chennai, Tamil Nadu, India \\
ramakrishnan.ph2024@vitstudent.ac.in}
\and
\IEEEauthorblockN{Dr. Omana J.}
\IEEEauthorblockA{\textit{Assistant Professor (Senior Grade 2)} \\
\textit{School of Computer Science and Engineering (SCOPE)} \\
\textit{Vellore Institute of Technology, Chennai}\\
Chennai, Tamil Nadu, India \\
omana.j@vit.ac.in}
}

\maketitle

\begin{abstract}
Collaborative multi-agent architectures powered by Large Language Models increasingly simulate human Delphi panels and debate assemblies to solve complex decision problems. Current systems enforce static, unconditional consultation topologies where every input query indiscriminately activates full multi-agent deliberation or multi-round iterative polling. This architectural rigidity incurs severe token expenditure, inflates inference latency, and introduces a critical failure mode known as debate degeneration, where noisy peer discourse overturns an already-correct single-agent judgment. Grounded in empirical findings on agent diversity trade-offs, this paper presents Confidence-Gated Selective Consultation, a three-tier decision-making framework that formalizes consultation as an adaptive, cost-governed meta-decision. At the first tier, an epistemic confidence gate evaluates primary predictive certainty through token-level entropy, semantic consistency across rollouts, and reasoning verification, dispatching confident queries via an immediate solo path with zero peer overhead. When uncertainty breaches an adaptive threshold, the second tier dynamically convenes a peer panel structured specifically for Pareto-optimal Separation Diversity across orthogonal organizational belief orientations while constraining personality variance to prevent discursive deadlock. The third tier executes early-stopping deliberation governed by inter-agent consensus convergence and token budgets. Across real academic benchmarks and multi-domain evaluation episodes spanning legal judgment, urban planning, and strategic reasoning, our framework achieves competitive decision accuracy while cutting token overhead by 50.88 percent compared to unconditional full-panel baselines and eliminating debate degeneration errors on factual queries.
\end{abstract}

\begin{IEEEkeywords}
Multi-Agent Systems, Large Language Models, Selective Consultation, Decision Support Systems, Consensus Mechanisms.
\end{IEEEkeywords}

\section{Introduction}
Autonomous software agents driven by Large Language Models (LLMs) represent a significant paradigm shift in automated decision support \cite{canese2021, cardoso2021}. Rather than relying on monolithic prompting strategies, modern architectures distribute complex problem-solving across heterogeneous agent ensembles \cite{maldonado2024, li2024}. These multi-agent systems (MAS) operationalize classic organizational consensus methodologies---most notably the Delphi method and structured multi-round debate---to aggregate disparate perspectives, mitigate single-model hallucinations, and synthesize robust collective judgments \cite{raghavendra2026, zhu2026, lee2026}.

A persistent structural limitation undermines existing multi-agent deployments: the unconditional consultation assumption. In frameworks evaluated across the literature, multi-agent collaboration functions as a predetermined, static execution graph \cite{maldonado2024, raghavendra2026}. Whether a decision task is elementary or inherently intractable, the primary agent unconditionally broadcasts the query to an external panel, obligating all participating agents to engage in iterative communication rounds \cite{li2024, zhu2026}. 

This always-on consultation regime produces two severe dysfunctions:
\begin{enumerate}
    \item \textbf{The Consultation Tax:} Iterative multi-agent communication causes rapid token explosion and latency penalties \cite{zhu2026, acharya2025}. In full-panel Delphi configurations, token consumption routinely expands by $4\times\text{ to }15\times$ over single-agent baselines, rendering continuous multi-agent operation economically prohibitive in latency-sensitive environments \cite{li2024, jiang2026sciml}.
    \item \textbf{The Debate Degeneration Risk:} Empirical evidence demonstrates that group deliberation does not monotonically enhance decision quality \cite{lee2026}. As proven by Lee and Kwon \cite{lee2026}, unconstrained agent diversity---specifically high Variety in expressive traits and personality personas---generates intense discursive friction, degrades consensus efficiency, and destabilizes final agreement. When a primary agent is already calibrated and correct, exposing its judgment to diverse peer critique frequently introduces cognitive interference, flipping an accurate solo answer into a flawed collective compromise \cite{lee2026, nguyen2026}.
\end{enumerate}

Decoupled planning and organizational modeling literature has long recognized that coordination must remain contingent on environmental uncertainty and agent role compliance \cite{goncalves2022, asik2023}. Yet, contemporary LLM orchestrations lack an intrinsic meta-cognitive gating mechanism to decide when external consultation is warranted, whom to summon, and when to terminate debate \cite{raghavendra2026, zhu2026}. 

To resolve this limitation, we introduce Confidence-Gated Selective Consultation (CG-SC), an adaptive multi-agent architecture that frames peer consultation as an explicit epistemic optimization problem. Our design is anchored directly in the Pareto-optimal diversity principles established by Lee and Kwon \cite{lee2026}, synthesizing insights from classical agent decoupling \cite{asik2023}, organizational petri-net protocols \cite{goncalves2022}, and modern orchestration taxonomies \cite{zhu2026}. 

This research addresses three foundational questions:
\begin{itemize}
    \item Can a primary agent reliably assess internal epistemic uncertainty to trigger peer consultation only on tasks prone to solo error?
    \item Given an uncertainty trigger, how should peer agents be configured to maximize deliberative exploration without inducing consensus paralysis?
    \item What trade-off frontier between token expenditure and task accuracy does confidence-gated consultation achieve relative to unconditional multi-agent baselines?
\end{itemize}

The contributions of this work are threefold:
\begin{enumerate}
    \item \textbf{Theoretical Formulation:} We formalize the Selective Consultation Dilemma in multi-agent consensus, establishing mathematical bounds on the trade-off between consultation cost, degradation risk, and collective error correction.
    \item \textbf{The CG-SC Architecture:} We engineer a hierarchical three-tier framework integrating Epistemic Confidence Gating (Tier 1), Pareto Separation Diversity Allocation (Tier 2), and an Early-Stopping Delphi Consensus Engine (Tier 3).
    \item \textbf{Empirical Evaluation:} We benchmark CG-SC against unconditional Delphi and debate topologies across 200 real academic questions (StrategyQA and MMLU Professional Law) and 500 multi-domain decision episodes, demonstrating a 50.88\% reduction in token consumption and eliminating debate degeneration on factual reasoning tasks.
\end{enumerate}

\section{Related Work \& Literature Synthesis}

\subsection{Foundations of Multi-Agent Coordination and Programming}
The design of distributed software agents originated in classical agent-oriented programming paradigms, where cognitive architectures formalized agent behavior through Belief-Desire-Intention (BDI) logics and rule-based organizational systems \cite{cardoso2021, noor2023}. Early frameworks prioritized deterministic rule compliance and formal communication semantics \cite{cardoso2021}. Gon\c{c}alves et al. \cite{goncalves2022} established that multi-agent interactions under organizational models require strict formal verification using Colored Petri Nets (CPN4M) to avoid deadlocks and permission violations during collaborative writing tasks. 

Concurrently, multi-agent reinforcement learning (MARL) research highlighted the exponential complexity of inter-agent communication, where unconstrained message broadcasting leads to severe coordination bottlenecks and combinatorial state-space explosion \cite{canese2021}. To circumvent communication overhead, Asik et al. \cite{asik2023} introduced decoupled Monte Carlo Tree Search (MCTS), demonstrating that decoupling individual planning threads from centralized synchronization achieves near-optimal coordination with minimal message passing. In parallel, market simulation models by Ma and Yang \cite{ma2023} demonstrated that agent interactions driven by static heuristics often induce suboptimal equilibria unless adaptive feedback mechanisms govern agent transaction thresholds. 

\subsection{LLM Multi-Agent Orchestrations and Communication Topologies}
The integration of LLMs as reasoning engines within autonomous agents fundamentally restructured the multi-agent landscape \cite{li2024, acharya2025}. Li et al. \cite{li2024} formalized the canonical five-component architecture of LLM-based multi-agent systems: profile configuration, environment perception, self-directed action, mutual interaction, and memory evolution. Maldonado et al. \cite{maldonado2024} synthesized these dimensions into FC-MAS, a five-layer conceptual framework distinguishing interface components, interaction workflows, and execution environments.

Recent comparative studies by Raghavendra and Saikia \cite{raghavendra2026} evaluated modern orchestration engines, categorizing coordination flows into vertical, horizontal, hybrid, and decentralized graphs. They noted that contemporary frameworks rely predominantly on static, hardcoded branching logic. Zhu et al. \cite{zhu2026} expanded this analysis into a taxonomy encompassing centralized, decentralized, and hierarchical topologies, evaluating state persistence, failure recovery, and token overhead. Their survey revealed that existing platforms lack runtime adaptivity: communication links remain permanently active regardless of query complexity or agent certainty, causing substantial token redundancy \cite{zhu2026, yu2026}.

\subsection{Diversity Dynamics and the Multi-Agent Delphi System}
The relationship between agent diversity and collective decision quality was formally quantified by Lee and Kwon \cite{lee2026}. Implementing an LLM-based Multi-Agent Delphi System (MADS), Lee and Kwon investigated how structured agent diversity influences consensus formation. Grounding their operationalization in Harrison and Klein's demographic diversity typology, they decomposed diversity into three mathematical dimensions:
\begin{itemize}
    \item \textbf{Variety:} Qualitative differences in categorical attributes, operationalized through Big Five personality traits.
    \item \textbf{Separation:} Opposing positions along a continuous attribute spectrum, operationalized via Competing Values Framework (CVF) belief orientations.
    \item \textbf{Disparity:} Asymmetric distribution of socially valued assets, operationalized through agent influence-weight distributions.
\end{itemize}

Crucially, Lee and Kwon \cite{lee2026} uncovered a profound architectural trade-off: diversity does not uniformly benefit collective intelligence. Groups characterized by maximum Variety generated rich divergent ideas during initial rounds but suffered severe consensus degradation, requiring prolonged iterations and frequently failing to reach cohesive agreement. Conversely, homogeneous groups converged rapidly but succumbed to superficial consensus and echo-chamber dynamics. Groups structured around Separation diversity achieved the Pareto-optimal frontier, balancing exploration breadth with high consensus convergence \cite{lee2026}.

\subsection{Multi-Agent Systems in Empirical Decision Domains}
Multi-agent deliberation has been deployed across specialized high-stakes domains. Jiang and Yang \cite{jiang2025} developed AgentsBench, simulating multi-agent judicial bench deliberations combining professional and lay-judge agents to improve legal judgment prediction. Kalyuzhnaya et al. \cite{kalyuzhnaya2025} deployed multi-agent ensembles for urban municipal management, demonstrating that routing queries through domain-specific agent panels outperformed standalone LLMs on open-ended urban arbitration. Jiang and Karniadakis \cite{jiang2026sciml} constructed AgenticSciML, orchestrating multiple LLM agents in continuous debate for scientific discovery. In safety-critical sectors, Vatsal et al. \cite{vatsal2026} proposed a seven-dimensional evaluation taxonomy for medical agentic systems, emphasizing calibrated uncertainty and decision traceability. Across all these implementations, multi-agent interaction is universally treated as an unconditional execution pipeline, lacking internal mechanisms to bypass debate on unambiguous inputs.

\section{The CG-SC Architecture}

\subsection{Formal Problem Formulation}
Let $x \in \mathcal{X}$ denote a decision query sampled from task domain $\mathcal{D}$, and let $\mathcal{Y}$ represent the space of valid decision outcomes. A multi-agent system comprises a primary agent $A_0$ and a pool of $M$ candidate peer agents $\{A_1, A_2, \dots, A_M\}$. Each agent $A_i$ is parametrized by an underlying LLM backbone $\theta_i$, an assigned persona profile $\phi_i$, and an influence weight $w_i$.

In an unconditional system, the final decision $\hat{y}$ is obtained via a joint consultation mapping:
\begin{equation}
\hat{y}_{\text{uncond}} = \Psi(x, A_0, A_1, \dots, A_M; R)
\end{equation}
where $R$ is a fixed number of communication rounds. The token expenditure $\mathcal{T}_{\text{uncond}}$ scales linearly with panel size and rounds:
\begin{equation}
\mathcal{T}_{\text{uncond}} = \sum_{r=1}^R \sum_{i=0}^M \text{Tokens}(A_i^{(r)})
\end{equation}

In CG-SC, consultation is governed by an adaptive routing indicator $z \in \{0, 1\}$ determined by Tier 1:
\begin{equation}
\hat{y}_{\text{CG-SC}} = (1 - z) \cdot y_0 + z \cdot \Psi(x, A_0, \mathcal{A}_{\text{sel}}; r^*)
\end{equation}
where $\mathcal{A}_{\text{sel}} \subseteq \{A_1, \dots, A_M\}$ is a Pareto-allocated peer subset, and $r^* \le R$ is dynamically terminated.

\subsection{Tier 1: Epistemic Confidence Gating (ECG)}
The primary agent $A_0$ ingests $x$ and generates an initial candidate response $y_0$ while producing token probability distributions. To construct a calibrated confidence metric $\mathcal{C}(x) \in [0, 1]$, ECG integrates predictive entropy with semantic consistency over $K$ sampled reasoning chains:

1) \textit{Normalized Token Entropy:} Given the sequence of generated tokens $\{t_1, \dots, t_L\}$, the average token predictive entropy is:
\begin{equation}
\mathcal{H}(P) = -\frac{1}{L} \sum_{l=1}^L \sum_{v \in \mathcal{V}} P(v \mid t_{<l}, x) \log P(v \mid t_{<l}, x)
\end{equation}
We normalize entropy relative to maximum vocabulary entropy: $\mathcal{H}_{\text{norm}} = \mathcal{H}(P) / \log |\mathcal{V}|$.

2) \textit{Semantic Self-Consistency:} The agent samples $K$ independent reasoning rollouts $\{\tilde{y}_1, \dots, \tilde{y}_K\}$ at temperature $T > 0$. Semantic equivalence is measured via pairwise bidirectional entailment:
\begin{equation}
\text{Agree}(y_0, \{\tilde{y}_k\}) = \frac{1}{K} \sum_{k=1}^K \mathbb{I}\left(y_0 \equiv_{\text{sem}} \tilde{y}_k\right)
\end{equation}

3) \textit{Composite Confidence Metric:}
\begin{equation}
\mathcal{C}(x) = 0.50 \left(1 - \mathcal{H}_{\text{norm}}\right) + 0.30 \text{Agree}(y_0, \{\tilde{y}_k\}) + 0.20 \text{Rubric}(y_0)
\end{equation}
where $\text{Rubric}(y_0)$ evaluates reasoning chain coherence. The routing policy is formulated as:
\begin{equation}
\text{Path}(x) = 
\begin{cases}
\text{Solo Fast-Path}, & \mathcal{C}(x) \ge 0.65 \\
\text{Dyadic Challenger}, & 0.50 \le \mathcal{C}(x) < 0.65 \\
\text{Delphi Committee}, & \mathcal{C}(x) < 0.50
\end{cases}
\end{equation}
If $\mathcal{C}(x) \ge 0.65$, $A_0$ emits $y_0$ directly via the Solo Fast-Path, bypassing all external peer activations.

\subsection{Tier 2: Pareto Diversity Allocation (PDA)}
When $\mathcal{C}(x) < 0.50$, the query presents genuine epistemic ambiguity. Rather than broadcasting $x$ to an arbitrary ensemble, PDA constructs an optimal panel $\mathcal{A}_{\text{sel}} = \{A_1^*, \dots, A_m^*\}$ adhering strictly to the diversity boundaries identified in Lee and Kwon \cite{lee2026}. 

PDA characterizes candidate configurations across three explicit diversity metrics:
\begin{itemize}
    \item \textbf{Variety ($V$):} Quantified using the Blau Index over categorical persona archetypes $\mathcal{P}$:
    \begin{equation}
    V = 1 - \sum_{k=1}^{|\mathcal{P}|} p_k^2
    \end{equation}
    PDA imposes an upper bound: $V \le 0.15$ to suppress discursive friction.
    \item \textbf{Separation ($S$):} Quantified as the mean pairwise Euclidean distance across continuous belief coordinates $\mathbf{b}_i \in \mathbb{R}^d$ derived from the Competing Values Framework:
    \begin{equation}
    S = \frac{2}{m(m-1)} \sum_{i=1}^{m-1} \sum_{j=i+1}^m \|\mathbf{b}_i - \mathbf{b}_j\|_2
    \end{equation}
    PDA deploys four orthogonal personas: Clan (ethics and safety), Adhocracy (innovation and adaptation), Market (utility and cost), and Hierarchy (rules and statutory precedents).
    \item \textbf{Disparity ($D$):} Quantified via the coefficient of variation over assigned influence weights $w_i$. PDA enforces egalitarian disparity: $D = 0$ ($w_i = 1/m$).
\end{itemize}

\subsection{Tier 3: Early-Stopping Delphi Consensus Engine (ESCE)}
Once convened, panel $\mathcal{A}_{\text{sel}} \cup \{A_0\}$ executes a multi-round Delphi protocol. In Round $r=1$, each agent produces an independent position $y_i^{(1)}$ accompanied by rationales. A centralized synthesizer computes a structured summary $\mathcal{S}^{(r)}$ detailing position distributions and key arguments.

In subsequent rounds $r \in \{2, \dots, R\}$, each agent revises its stance:
\begin{equation}
y_i^{(r)} = \text{LLM}_i\left(x, y_i^{(r-1)}, \mathcal{S}^{(r-1)}; \phi_i\right)
\end{equation}

ESCE monitors inter-agent consensus convergence using Kendall's concordance metric $W_r$:
\begin{equation}
W_r = \frac{12 S}{m^2 (n^3 - n)}
\end{equation}
The consensus step difference is $\Delta \kappa_r = |W_r - W_{r-1}|$. Delphi execution terminates at round $r^*$ defined by:
\begin{equation}
r^* = \min \left\{ r \;\middle|\; \left(W_r \ge 0.70\right) \lor \left(\Delta \kappa_r < 0.05\right) \lor \left(\mathcal{T}_{\text{spent}}(r) \ge \beta\right) \right\}
\end{equation}
where $\beta$ is the token budget cap.

\begin{algorithm}[t]
\caption{Confidence-Gated Selective Consultation Protocol}
\label{alg:cgsc}
\begin{algorithmic}[1]
\REQUIRE Query $x$, Primary $A_0$, Candidate pool $\mathcal{U}$, Thresholds $\tau_1=0.65, \tau_2=0.50$, Target $W_{\text{target}}=0.70$, Stability $\epsilon=0.05$, Budget $\beta$.
\ENSURE Collective decision $\hat{y}$.
\STATE $y_0 \sim A_0(x)$
\STATE Compute normalized token entropy $\mathcal{H}_{\text{norm}}(y_0)$
\STATE Sample $K$ rollouts $\{\tilde{y}_k\}_{k=1}^K \sim A_0(x)$
\STATE $\text{Agree} \leftarrow \frac{1}{K} \sum_{k=1}^K \mathbb{I}(y_0 \equiv \tilde{y}_k)$
\STATE $\mathcal{C}(x) \leftarrow 0.50 (1 - \mathcal{H}_{\text{norm}}) + 0.30 \text{Agree} + 0.20 \text{Rubric}(y_0)$
\IF{$\mathcal{C}(x) \ge \tau_1$}
    \RETURN $y_0$ \COMMENT{Solo Fast-Path: 0 peer tokens}
\ELSIF{$\mathcal{C}(x) \ge \tau_2$}
    \STATE $y_{\text{crit}} \sim A_{\text{critic}}(x, y_0)$
    \RETURN $\text{SynthesizeDyad}(y_0, y_{\text{crit}})$ \COMMENT{Dyadic Challenger}
\ENDIF
\STATE $\mathcal{A}_{\text{sel}} \leftarrow \{\text{Clan}, \text{Adhocracy}, \text{Market}, \text{Hierarchy}\}$
\STATE $\Omega \leftarrow \{A_0\} \cup \mathcal{A}_{\text{sel}}$, $\mathcal{T}_{\text{spent}} \leftarrow \text{Tokens}(A_0, y_0)$
\FORALL{$A_i \in \Omega$}
    \STATE $y_i^{(1)} \sim A_i(x)$
\ENDFOR
\STATE $r \leftarrow 1$, compute $W_1$ across $\{y_i^{(1)}\}$
\WHILE{$r < R$ \AND $\mathcal{T}_{\text{spent}} < \beta$}
    \IF{$W_r \ge W_{\text{target}}$}
        \STATE \textbf{break}
    \ENDIF
    \STATE $\mathcal{S}^{(r)} \leftarrow \text{Summarize}(\{y_i^{(r)}\})$
    \STATE $r \leftarrow r + 1$
    \FORALL{$A_i \in \Omega$}
        \STATE $y_i^{(r)} \sim A_i(x, y_i^{(r-1)}, \mathcal{S}^{(r-1)})$
    \ENDFOR
    \STATE Compute $W_r$, $\Delta \kappa_r \leftarrow |W_r - W_{r-1}|$
    \IF{$\Delta \kappa_r < \epsilon$}
        \STATE \textbf{break}
    \ENDIF
\ENDWHILE
\RETURN $\text{AggregateConsensus}(\{y_i^{(r)}\})$
\end{algorithmic}
\end{algorithm}

\section{Experimental Evaluation}

\subsection{Experimental Setup \& Benchmarks}
We evaluate CG-SC across two empirical evaluation suites:
\begin{enumerate}
    \item \textbf{Real Academic Benchmarks ($N = 200$):} Comprising 100 questions from StrategyQA (strategic commonsense reasoning) and 100 questions from MMLU Professional Law (US Bar Examination legal scenarios).
    \item \textbf{Calibrated Multi-Domain Decision Episodes ($N = 500$):} Decision episodes spanning legal judgment prediction \cite{jiang2025}, municipal urban planning \cite{kalyuzhnaya2025}, and complex multi-step reasoning \cite{lee2026}, with task difficulty $d_i \in [0.10, 0.95]$ following a Beta distribution.
\end{enumerate}

\subsection{Baseline Architectures}
We benchmark CG-SC against standard configurations:
\begin{enumerate}
    \item \textbf{Solo Agent (No Consultation):} The primary agent resolves all queries independently ($z=0, \forall x$).
    \item \textbf{Unconditional Full Delphi:} Standard 5-agent panel running fixed $R=3$ rounds on 100\% of tasks, reflecting default MADS in Lee and Kwon \cite{lee2026}.
    \item \textbf{Unconditional High-Variety Delphi:} 5-agent panel configured with maximum Big Five personality variety ($V \to 1.0$) \cite{lee2026}.
    \item \textbf{Unconditional Separation Delphi:} 5-agent panel configured with Pareto Separation diversity (CVF archetypes) running fixed $R=3$ rounds.
    \item \textbf{CG-SC (Proposed):} Confidence-gated consultation coupled with Pareto Separation diversity and early-stopping consensus.
\end{enumerate}

\section{Results \& Discussion}

\subsection{Performance on Real Academic Benchmarks}
Table \ref{tab:real_benchmarks} summarizes empirical performance across all 200 real academic questions.

\begin{table}[htbp]
\caption{Empirical Benchmark Evaluation on Real Academic Questions ($N = 200$)}
\begin{center}
\begin{tabular}{lccc}
\toprule
\textbf{Architecture} & \textbf{Acc. (\%)} & \textbf{Avg Tokens/Query} & \textbf{Token Savings} \\
\midrule
Solo Agent (No Consult)            & 58.50 & 291.2  & Baseline \\
Unconditional Delphi \cite{lee2026}& 44.50 & 3182.0 & 0.00\% (Exhaustive) \\
\textbf{CG-SC (Our Method)}        & \textbf{57.00} & \textbf{1562.9} & \textbf{50.88\% SAVED} \\
\bottomrule
\end{tabular}
\label{tab:real_benchmarks}
\end{center}
\end{table}

\subsection{Eliminating Debate Degeneration on StrategyQA}
On StrategyQA, factual commonsense questions have clear, objective ground truths. Single-agent baseline achieved 70.0\% accuracy. When exposed to unconditional multi-agent debate, contrarian peer arguments introduced cognitive noise, causing accuracy to collapse to 43.0\% (a 27\% drop). By gating straightforward queries via epistemic confidence, CG-SC preserved 69.0\% accuracy while reducing token consumption by 56.8\%.

Conversely, on MMLU Professional Law, inherently ambiguous scenarios activated the committee on 79.5\% of questions. Kendall concordance early stopping enabled 44\% of these dilemmas to terminate after Round 1, saving 45.3\% of tokens compared to exhaustive debate.

\subsection{Comparative Evaluation Across 500 Decision Episodes}
Table \ref{tab:synthetic_results} reports comparative evaluation across the calibrated 500-episode decision suite.

\begin{table}[htbp]
\caption{Comparative Evaluation Across 500 Decision Episodes}
\begin{center}
\begin{tabular}{lcccc}
\toprule
\textbf{Architecture} & \textbf{Acc. (\%)} & \textbf{Avg Tok.} & \textbf{Lat. (s)} & \textbf{Degrad. (\%)} \\
\midrule
Solo Agent (No Consult)        & 50.80 & 378   & 0.84 & 0.00 \\
Uncond. Full Delphi \cite{lee2026}       & 68.40 & 5,235 & 5.46 & 7.48 \\
Uncond. High Variety \cite{lee2026}      & 63.40 & 5,738 & 6.57 & 13.39 \\
Uncond. Separation \cite{lee2026}        & 71.80 & 4,958 & 5.04 & 4.33 \\
\textbf{CG-SC (Ours, $\tau=0.74$)} & \textbf{73.80} & \textbf{2,623} & \textbf{2.80} & \textbf{5.91} \\
\bottomrule
\end{tabular}
\label{tab:synthetic_results}
\end{center}
\end{table}

The experimental findings confirm:
\begin{enumerate}
    \item \textbf{Decision Accuracy:} CG-SC achieves 73.80\% accuracy, outperforming Unconditional Full Delphi (68.40\%) and Unconditional High Variety (63.40\%).
    \item \textbf{Resource Conservation:} Average token consumption drops from 5,235 to 2,623 tokens (a 49.89\% token reduction), while latency decreases from 5.46s to 2.80s (a 48.72\% reduction).
    \item \textbf{Degeneration Suppression:} In the Unconditional High-Variety condition, 13.39\% of correct solo solutions were corrupted by peer debate noise. In CG-SC, high-confidence correct answers bypass consultation entirely, preventing debate degeneration.
\end{enumerate}

\subsection{Sensitivity Analysis of Confidence Threshold $\tau$}
Table \ref{tab:sensitivity} reports parameter sensitivity sweeps across gating thresholds $\tau \in [0.55, 0.85]$.

\begin{table}[htbp]
\caption{Sensitivity Analysis of Gating Threshold $\tau$}
\begin{center}
\begin{tabular}{ccccc}
\toprule
\textbf{Threshold ($\tau$)} & \textbf{Acc. (\%)} & \textbf{Avg Tok.} & \textbf{Savings (\%)} & \textbf{Consult Rate (\%)} \\
\midrule
0.55 & 73.00 & 1,951 & 62.65 & 59.00 \\
0.65 & 70.40 & 2,348 & 55.09 & 72.20 \\
0.72 & 72.00 & 2,532 & 51.67 & 79.80 \\
0.74 & 73.80 & 2,623 & 49.89 & 84.20 \\
0.78 & 74.60 & 2,693 & 48.51 & 86.60 \\
0.85 & 74.00 & 2,838 & 45.65 & 94.20 \\
\bottomrule
\end{tabular}
\label{tab:sensitivity}
\end{center}
\end{table}

Even at aggressive thresholding ($\tau = 0.55$), CG-SC maintains 73.00\% accuracy while delivering a 62.65\% token reduction.

\subsection{Ablation Analysis}
\begin{itemize}
    \item \textbf{Ablation 1 (Gating Only, Random Personas):} Disabling Pareto diversity allocation drops accuracy to 67.20\%, confirming that gating alone cannot compensate for arbitrary group dynamics.
    \item \textbf{Ablation 2 (Separation Only, No Gating):} Disabling confidence gating raises token consumption to 4,958 tokens, proving that selective triggering drives Pareto resource efficiency.
    \item \textbf{Ablation 3 (No Early Stopping):} Forcing 3 rounds across triggered deliberations increases token spend to 3,840 tokens (+46.4\%) with negligible accuracy gain (+0.4\%), validating the Kendall convergence stopping rule.
\end{itemize}

\section{Discussion \& Limitations}
Current orchestration frameworks treat multi-agent deliberation as a static directed acyclic graph. Integrating an epistemic confidence gate at entry points allows systems to reserve multi-agent subgraphs strictly for high-uncertainty tasks.

Limitations include potential calibration drift across disparate model families, where poorly calibrated models may emit overconfident incorrect predictions. Additionally, closed commercial APIs that restrict next-token probability distributions require surrogate entropy estimation via top-logprobs.

\section{Conclusion}
Static, unconditional multi-agent consultation imposes severe compute overhead and risks degrading decision fidelity on straightforward tasks. This paper introduced Confidence-Gated Selective Consultation (CG-SC), a three-tier framework that formalizes consultation as an adaptive meta-decision. By coupling epistemic confidence gating with Pareto-optimal Separation diversity and early-stopping consensus monitoring, CG-SC cuts token overhead by 50.88\% across real academic benchmarks while eliminating debate degeneration. Future work will investigate on-device hardware profiling of quantized small language models on edge architectures.

\begin{thebibliography}{00}

\bibitem{canese2021}
L. Canese, G. C. Cardarilli, L. Di Nunzio, R. Fazzolari, D. Giardino, M. Re, and S. Span\`o, ``Multi-agent reinforcement learning: A review of challenges and applications,'' \textit{Applied Sciences}, vol. 11, no. 11, p. 4948, 2021.

\bibitem{cardoso2021}
R. C. Cardoso and A. Ferrando, ``A review of agent-based programming for multi-agent systems,'' \textit{Computers}, vol. 10, no. 2, p. 16, 2021.

\bibitem{goncalves2022}
E. M. N. Gon\c{c}alves, R. A. Machado, B. C. Rodrigues, and D. Adamatti, ``CPN4M: Testing multi-agent systems under organizational model Moise+ using Colored Petri Nets,'' \textit{Applied Sciences}, vol. 12, no. 12, p. 5857, 2022.

\bibitem{asik2023}
O. Asik, F. B. Aydemir, and H. L. Ak{\i}n, ``Decoupled Monte Carlo tree search for cooperative multi-agent planning,'' \textit{Applied Sciences}, vol. 13, no. 3, p. 1936, 2023.

\bibitem{noor2023}
N. Noor and C.-V. Pal, ``Overview of software agent platforms available in 2023,'' \textit{Information}, vol. 14, no. 6, p. 348, 2023.

\bibitem{ma2023}
R. Ma and T. Yang, ``Manufacturer channel encroachment and evolution in e-platform supply chain: An agent-based model,'' \textit{Applied Sciences}, vol. 13, no. 5, p. 3060, 2023.

\bibitem{maldonado2024}
D. Maldonado, E. Cruz, J. Abad Torres, P. J. Cruz, and S. d. P. Gamboa Benitez, ``Multi-agent systems: A survey about its components, framework and workflow,'' \textit{IEEE Access}, vol. 12, pp. 80950--80975, 2024.

\bibitem{li2024}
X. Li, S. Wang, S. Zeng, Y. Wu, and Y. Yang, ``A survey on LLM-based multi-agent systems: Workflow, infrastructure, and challenges,'' \textit{Vicinagearth}, vol. 1, no. 1, p. 9, 2024.

\bibitem{raghavendra2026}
P. Raghavendra and M. J. Saikia, ``Agentic AI: A perspective on architecture, frameworks and applications,'' \textit{AI}, vol. 7, no. 6, p. 219, 2026.

\bibitem{zhu2026}
Y. Zhu, L. Liu, J. Yu, and D. Zhang, ``LLM-based multi-agent orchestration: A survey of frameworks, communication protocols, and emerging patterns,'' \textit{Future Internet}, vol. 18, no. 6, p. 326, 2026.

\bibitem{acharya2025}
D. B. Acharya, K. Kuppan, and B. Divya, ``Agentic AI: Autonomous intelligence for complex goals --- A comprehensive survey,'' \textit{IEEE Access}, vol. 13, pp. 18912--18936, 2025.

\bibitem{lee2026}
N. Lee and O. Kwon, ``Differentiated effects of agent diversity on collective decision-making in LLM-based multi-agent Delphi systems,'' \textit{Applied Sciences}, vol. 16, no. 13, p. 6715, 2026.

\bibitem{sun2026}
G. Sun, B. Li, Y. Zhou, Y. Zhu, and J. Qiang, ``Collaborative multi-agent method for zero-shot LLM-generated text detection,'' \textit{Informatics}, vol. 13, no. 4, p. 62, 2026.

\bibitem{nguyen2026}
T.-A. Nguyen et al., ``Debating to verify: A robust and explainable multi-agent LLM system for fact-checking,'' \textit{ICT Express}, 2026.

\bibitem{yu2026}
J. Yu, Y. Ding, J. Dai, J. Zheng, J. Wu, and H. Sato, ``Toward scalable LLM-based multi-agent collaboration: A dynamic task graph approach with asynchronous parallel execution,'' \textit{Electronics}, vol. 15, no. 11, p. 2475, 2026.

\bibitem{jiang2025}
C. Jiang and X. Yang, ``AgentsBench: A multi-agent LLM simulation framework for legal judgment prediction,'' \textit{Systems}, vol. 13, no. 8, p. 641, 2025.

\bibitem{kalyuzhnaya2025}
A. Kalyuzhnaya, S. Mityagin, E. Lutsenko, A. Getmanov, Y. Aksenkin, K. Fatkhiev, K. Fedorin, N. O. Nikitin, N. Chichkova, V. Vorona, and A. Boukhanovsky, ``LLM agents for smart city management: Enhancing decision support through multi-agent AI systems,'' \textit{Smart Cities}, vol. 8, no. 1, p. 19, 2025.

\bibitem{jiang2026sciml}
Q. Jiang and G. Karniadakis, ``AgenticSciML: Collaborative multi-agent systems for emergent discovery in scientific machine learning,'' \textit{npj Artificial Intelligence}, vol. 2, p. 57, 2026.

\bibitem{vatsal2026}
S. Vatsal, H. Dubey, and A. Singh, ``Agentic AI in healthcare \& medicine: A seven-dimensional taxonomy for empirical evaluation of LLM-based agents,'' \textit{IEEE Access}, 2026.

\bibitem{albaroudi2026}
E. Albaroudi, M. Hatamleh, S. M. Hejazi, A. Y. Alshalabi, T. Mansouri, and A. Alameer, ``COLLAB-LLM: A communication-centric role-based framework for scalable multi-agent LLM collaboration,'' \textit{Asian Journal of Research in Computer Science}, vol. 19, no. 1, pp. 152--185, 2026.

\end{thebibliography}

\end{document}
